\chapter{Деньги между банками}\label{ch:banks}

Клиент банка~А переводит сто тысяч рублей клиенту банка~Б;
в~мобильном приложении это одно действие: ввести сумму, нажать
кнопку, дождаться статуса \enquote{исполнено}. На~каждой ступени
этого пути меняются форма денег, обязанная сторона и~система
учёта, в~которой это обязательство зафиксировано. Что именно
движется между участниками, кто имеет право менять учётные записи
и~в~какой момент перевод считается исполненным~--- три вопроса
этой главы.

\section{Формы денег}\label{sec:banks-forms-of-money}

Деньги существуют в~нескольких формах, и~каждая подчиняется
собственным правилам учёта, хранения и~передачи. Разработчик,
впервые столкнувшись с~платёжной системой, обычно видит деньги
как поле \texttt{amount} с~типом \texttt{decimal}, которое можно
увеличить или уменьшить. Эта модель работает в~коде и~перестаёт
работать в~тот момент, когда деньги переходят из~одной формы
в~другую.

Закон различает три формы рубля: наличные, безналичные деньги
и, с~2023 года, цифровой рубль. Отдельно 161-ФЗ вводит
понятие \textit{электронных денежных средств}~--- юридически это
особая конструкция, стоящая рядом с~наличной и~безналичной
формами~\cite{fz-161}.

\subsection*{Наличные}

Передача наличных не требует посредника и~завершается в~момент
физического вручения. Покупатель протягивает купюру продавцу,
продавец кладёт её в~кассу; расчёт окончен, никакой банк не
участвует, никакая запись в~реестре не нужна. Вопрос \enquote{дошли
ли деньги?} не возникает в~принципе: деньги перешли из~рук в~руки
на~глазах у~обоих участников.

\highlight{Юридически наличные~--- обязательство Центрального банка}.
Каждая банкнота свидетельствует о~том, что ЦБ должен её предъявителю
определённую сумму, и~это обязательство обеспечено всеми активами
Банка России~\cite{fz-86}. Эмиссия наличных~--- монопольное право ЦБ;
безналичные деньги, напротив, создают коммерческие банки, выдавая кредиты.

\subsection*{Безналичные деньги}

Когда клиент вносит сто тысяч рублей на~счёт в~банке, наличные
перестают существовать как наличные и~превращаются в~запись
на~банковском счёте, то~есть в~обязательство коммерческого банка
перед клиентом. У~клиента больше нет той же купюры~--- у~него
появилось требование к~банку, юридически отличное от~наличных,
с~которыми он только что пришёл.

Различие между наличными и~безналичными в~обыденной жизни незаметно,
пока банк выдаёт наличные по~первому требованию. Базовый механизм страхования вкладов покрывает
до~1,4~миллиона рублей на~одного вкладчика в~одном банке,\footnote{С~30~октября
2025~года для рублёвых вкладов, удостоверенных безотзывными сберегательными
сертификатами сроком свыше трёх лет, действует отдельное страховое покрытие
до~2,8~миллиона рублей; оно считается отдельно от~базового лимита, общий потолок
по~двум видам вкладов достигает 4,2~миллиона рублей.} и~существует именно
потому, что \highlight{безналичные деньги~--- обязательства коммерческого
банка}, а~коммерческий банк может обанкротиться. Банкнота в~кармане
от~банкротства конкретного банка не зависит~\cite{fz-177}.

Для платёжных систем существенно другое свойство безналичных
денег: их передача требует изменения записей минимум в~двух
местах. Если оба счёта находятся в~одном банке, операция сводится
к~внутрибанковской проводке и~выполняется мгновенно. Банк контролирует
оба счёта и~меняет записи атомарно. Если счёта расположены в~разных
банках, две независимые учётные системы должны \textit{согласованно}
изменить свои записи, не потеряв и~не удвоив деньги в~процессе.
Для этого и~существуют \textit{корреспондентские счета}~--- счета,
которые один банк открывает в~другом банке для проведения межбанковских
расчётов.

\subsection*{Электронные денежные средства и~цифровой рубль}

161-ФЗ вводит ещё одну форму: электронные денежные средства (ЭДС),
определяя их как денежные средства, предварительно предоставленные
лицом оператору, который учитывает их без открытия банковского счёта
и~обязуется переводить их по~распоряжениям этого лица~\cite{fz-161}.
ЭДС реализуются как кошельки вроде тех, что когда-то
предлагали Яндекс.Деньги (ныне ЮMoney) и~Qiwi: клиент пополняет
кошелёк, оператор хранит баланс и~по команде переводит средства
получателю. Банковский счёт не открывается, средства на~кошельке
не покрываются страхованием вкладов.

Для инженера, проектирующего платёжный продукт, различие между
безналичными деньгами и~ЭДС проявляется в~лицензировании и~лимитах.
Оператор ЭДС обязан быть кредитной организацией (банком или
небанковской кредитной организацией, НКО), а~лимиты на~остаток
и~оборот зависят от~\textit{уровня идентификации клиента}: для
неперсонифицированных и~упрощённо идентифицированных пользователей
закон задаёт разные пороги по~остатку и~операциям, для
персонифицированного режима ограничения заметно мягче~\cite{fz-161}.

В~2023 году к~наличным и~безналичным деньгам добавилась третья форма
рубля: \highlight{цифровой рубль~--- прямое обязательство
Банка России}, хранящееся на~платформе ЦБ; коммерческие банки лишь
предоставляют клиентам доступ к~этой платформе~\cite{fz-340}.
По~обязанной стороне цифровой рубль ближе к~наличным, чем
к~безналичным; передача всегда проходит через платформу и~оставляет
прослеживаемый след. Архитектура и~взаимодействие с~СБП (системой
быстрых платежей, см.~главу~\ref{ch:sbp}) разобраны
в~главе~\ref{ch:cbdc}.

\section{Корреспондентские счета}\label{sec:banks-correspondent}

Банк~А списывает средства со~счёта отправителя~--- это его
внутренняя операция. Зачислить деньги на~счёт получателя в~банке~Б
он не может, потому что чужой учётной системой не управляет.
Решение этой проблемы появилось задолго до~компьютеров и~карточных
сетей~--- корреспондентские счета.

У~такого счёта две стороны, и~названия их закрепились ещё в~итальянской
банковской практике. Счёт, открытый \emph{нашим} банком \emph{в~чужом},
называется \textit{ностро-счётом} (от~итальянского \emph{nostro}~---
\enquote{наш}); тот же самый счёт с~точки зрения банка, в~котором он
открыт, называется \textit{лоро-счётом} (от~итальянского \emph{loro}~---
\enquote{их}). Один и~тот же счёт всегда одновременно ностро для одного
банка и~лоро для другого~--- асимметрия точки зрения, а~не два разных
счёта.

Пусть банк~А открыл корреспондентский счёт в~банке~Б (это ностро-счёт
банка~А и~лоро-счёт для банка~Б), и~остаток счёта покрывает перевод.
Клиент банка~А инициирует перевод. Банк~А списывает средства со~счёта
клиента и~отправляет банку~Б сообщение: \enquote{спиши сто тысяч
с~нашего ностро-счёта у~тебя и~зачисли на~счёт такого-то получателя}.
Банк~Б дебетует лоро-счёт банка~А, уменьшая остаток, который банк~А
держит у~него, и~кредитует счёт получателя. Перевод завершён.

Эта схема требует прямых двусторонних связей между каждой парой
банков: один из~них должен открыть счёт у~другого, а~чаще оба
открывают счёта друг у~друга. При тысяче банков в~стране число
возможных пар составляет около полумиллиона, и~поддерживать полмиллиона
двусторонних отношений нецелесообразно ни экономически, ни
операционно.

Сеть схлопывается в~звезду, когда появляется центральный корреспондент,
у~которого остальные банки держат корреспондентские счёта. В~России
эту роль выполняет Банк России. Банки переводят
средства через счета, открытые в~Банке России, и~через ту же систему
завершается расчёт по~карточным операциям внутри страны~\cite{cbr-ps}.
Когда клиент банка~А переводит деньги клиенту банка~Б, последовательность
событий такова:

\begin{itemize}
  \item Банк~А списывает средства со~счёта отправителя.
  \item Банк~А направляет платёжное поручение в~платёжную систему Банка России.
  \item Банк России дебетует корреспондентский счёт банка~А и~кредитует корреспондентский счёт банка~Б.
    Здесь происходит фактическое перемещение денег между банками. Оба счёта находятся в~одной учётной системе
    Центрального банка, и~он изменяет обе записи атомарно.
  \item Банк~Б зачисляет средства на~счёт получателя. Этот шаг отражает уже состоявшийся расчёт во~внутреннем учёте банка~Б.
\end{itemize}

\practicenote{%
  Поля \enquote{корр.\,счёт} и~\enquote{БИК} в~реквизитах для перевода~---
  именно эта инфраструктура. БИК (банковский идентификационный код)
  однозначно идентифицирует банк в~платёжной системе Банка России.
  Корреспондентский счёт~--- тот самый счёт, через который банк
  рассчитывается с~другими банками при посредничестве ЦБ.

  БИК относится к~межбанковским расчётам, тогда как BIN (Bank
  Identification Number)~--- к~карточным; оба идентифицируют банк,
  но в~разных контекстах.
}

В~международных расчётах единого центрального корреспондента нет.
Роль посредников берут на~себя крупные \textit{банки-корреспонденты},
и~цепочка растягивается до~трёх-четырёх звеньев: от~банка отправителя
через одного или нескольких посредников к~банку получателя. Каждое
звено образует пару корреспондентских счетов, каждая проводка означает
дебет одного счёта и~кредит другого. На~каждом переходе банк-посредник
может удержать комиссию, конвертировать валюту и~задержать перевод
для комплаенс-проверки на~соответствие регуляторным требованиям,
в~том числе санкционным и~антиотмывочным. Международный
банковский перевод дорог и~медлителен; карточные сети с~единой
расчётной инфраструктурой и~предсказуемым сроком расчёта стали для
бизнеса стандартом.

\section{Платёжная инфраструктура: ПС Банка России, SWIFT, СПФС}\label{sec:banks-infrastructure}

Корреспондентские счета определяют, \emph{где} учитываются деньги.
Они не отвечают на~вопрос, \emph{как} банк~А передаёт банку~Б (или
Центральному банку) распоряжение о~переводе, в~каком формате и~с~какими
гарантиями доставки. За~это отвечает платёжная инфраструктура~---
системы и~протоколы, через которые передаются платёжные сообщения
и~проводится расчёт.

\subsection*{Платёжная система Банка России}

Основу внутрироссийских межбанковских расчётов составляет платёжная
система Банка России. Через неё проходят и~обычные клиентские переводы,
и~налоговые платежи, и~операции на~межбанковском рынке~\cite{cbr-ps}.
Она работает в~двух режимах: рейсовом для рутинного потока и~срочном
(RTGS) для крупных межбанковских операций.

Обычные переводы накапливаются в~течение операционного дня
и~проходят несколько \textit{расчётных рейсов}. К~очередному рейсу банк
получает результат; окончательный расчёт включает проверку достаточности
средств на~корреспондентском счёте, формирование внутридневной очереди
и, где возможно, взаимозачёт. Этот режим покрывает основной объём
рутинных платежей: зарплаты, оплату поставщикам, коммунальные. Высокий
остаток на~корреспондентском счёте банку не нужен, потому что часть входящих
и~исходящих платежей зачитывается внутри рейса.

Срочные платежи проходят через сервис срочного перевода ПС Банка России~---
режим RTGS (Real-Time Gross Settlement, валовой расчёт в~реальном времени).
Каждый платёж проводится отдельно и~немедленно, без накопления и~без
взаимозачёта. Средства списываются с~корреспондентского счёта
банка-отправителя и~зачисляются на~счёт банка-получателя по~каждому
платежу отдельно~\cite{cbr-ps}. Срочный перевод стоит банку дороже,
чем несрочный рейсовый, и~даёт почти мгновенную окончательность расчёта.
Для крупных межбанковских сделок и~операций денежного рынка это решающее
преимущество.\footnote{Исторически
функцию RTGS внутри ПС Банка России выполняла отдельная компонента
БЭСП (банковские электронные срочные платежи), выведенная
из~эксплуатации 2~июля 2018~года; с~этого момента её
функционал перенесён в~Перспективные платёжные
сервисы (ППС) Банка России.}

\importantnote{%
  Сервис срочного перевода ПС Банка России и~система быстрых
  платежей (СБП) внешне похожи скоростью завершения, но решают
  разные задачи. Сервис срочного перевода работает
  как \textit{межбанковская} система RTGS для крупных и~срочных
  платежей между банками, тогда как СБП~--- \textit{розничная}
  платёжная система для мгновенных переводов физических лиц
  и~оплаты покупок, работающая через инфраструктуру НСПК
  (Национальной системы платёжных карт).
}

\subsection*{SWIFT и~СПФС}

Международные межбанковские расчёты традиционно работают через SWIFT
(Society for Worldwide Interbank Financial Telecommunication)~---
кооперативную организацию, основанную в~1973 году и~объединяющую
порядка 11\,500 участников более чем в~200 странах: банки, финансовые
организации и~крупные корпорации~\cite{swift-overview}.

SWIFT доставляет структурированное сообщение от~одного банка к~другому~---
например, инструкцию \enquote{переведи столько-то с~нашего ностро-счёта
у~тебя на~счёт такого-то клиента}. Фактическое движение происходит
по~корреспондентским счётам, описанным в~предыдущем разделе,
а~выполнение инструкции остаётся делом банка-получателя.

Сообщения SWIFT исторически передавались в~формате MT (Message Type):
MT103~--- клиентский перевод, MT202~--- межбанковский перевод, MT940~---
выписка по~счёту и~так далее. Параллельно отрасль переходит
к~сообщениям MX (XML-формат сообщений на~базе ISO~20022). К~22 ноября
2025 года SWIFT по~программе CBPR+ (Cross-Border Payments and Reporting Plus,
миграция трансграничных платежей и~отчётности на~ISO~20022) завершил период
сосуществования MT и~ISO~20022; для платёжных
инструкций ISO~20022 стал базовым форматом~\cite{swift-iso20022-milestone,swift-iso20022-implementation}.

СПФС (Система передачи финансовых сообщений)~--- инфраструктура обмена
финансовыми сообщениями Банка России, созданная как российская
альтернатива SWIFT. Банк России описывает СПФС как канал передачи
электронных сообщений по~финансовым операциям; сами расчёты кредитные
организации проводят по~корреспондентским счетам~\cite{cbr-spfs}.
\highlight{По~функциональности СПФС аналогична SWIFT: передаёт
структурированные финансовые сообщения и~параллельно поддерживает
несколько форматов}~--- SWIFT MT, ISO~20022 (MX), УФЭБС
(унифицированные форматы электронных банковских сообщений Банка
России) и~пользовательские форматы, согласованные участниками.

СПФС находится под контролем Банка России, данные хранятся
в~Российской Федерации, подключение регулируется российским
законодательством~--- три отличия от~SWIFT, ради которых система
и~строилась.

Банку, обслуживающему и~внутрироссийские, и~международные переводы,
приходится держать два канала: СПФС для внутренних операций и~SWIFT
(или его альтернативы)~--- для международных. Форматы сообщений
пересекаются (ISO~20022, SWIFT MT), но
\highlight{маршрутизация, правила идентификации участников и~требования
комплаенса различаются}, и~унифицировать их в~одном
модуле без потери корректности сложно.

\needspace{4\baselineskip}
\section{Банковский перевод и~карточная транзакция}\label{sec:banks-transfer-vs-card}

Карточная транзакция устроена сложнее классического банковского
перевода. Когда покупатель прикладывает карту к~терминалу, деньги
\emph{не переходят} к~продавцу сразу. Банк, выпустивший карту
(эмитент), лишь подтверждает готовность заплатить; банк продавца
(эквайер) получит деньги позже, а~фактическое межбанковское
урегулирование произойдёт после клиринга и~расчёта (роли участников
описаны в~разделе~\ref{sec:ecosystem-participants}).

\textit{Авторизация} занимает доли секунды.
Терминал отправляет запрос по~платёжной инфраструктуре, эмитент
проверяет карту, баланс, лимиты и~антифрод
и~возвращает код \texttt{00} (стандартный код одобрения по~ISO~8583,
см.~главу~\ref{ch:iso8583}) или иной код отказа за~время, которое
покупатель едва ли замечает. При одобрении банк резервирует сумму
на~счёте держателя карты, но не списывает её окончательно
и~не переводит деньги в~расчёты.

\textit{Клиринг} (clearing) происходит позже, в~конце операционного дня. Эквайер собирает все
авторизованные операции за~день и~передаёт их в~карточную сеть
клиринговым файлом~--- структурированным набором записей с~реквизитами
транзакции, суммой, валютой и~ссылкой на~авторизацию. Карточная сеть
сопоставляет клиринговые записи с~авторизациями, рассчитывает комиссии
и~формирует итоговые расчётные позиции для банков-участников.

\textit{Расчёт} (settlement) завершает цикл. Карточная сеть или
уполномоченный расчётный банк инициирует фактическое перемещение
денег между корреспондентскими счётами банков. Эквайер
получает причитающуюся сумму за~вычетом комиссий, эмитент
окончательно списывает средства со~счёта
держателя карты. \highlight{Для внутрироссийских карточных операций
обработка и~клиринг проходят через инфраструктуру НСПК, а~завершение
расчётов происходит в~платёжной системе Банка России}~\cite{cbr-ps}.

\importantnote{%
  Разрыв между авторизацией и~расчётом обходится отрасли тремя
  устойчивыми эффектами:
  \begin{enumerate}
    \item холды (\textit{hold}) могут сохраняться по~нескольку дней;
    \item легитимна ситуация, при которой авторизация прошла на~одну сумму,
      а~в~клиринге пришла другая (типичный случай: ресторан, где чаевые
      добавляются после авторизации);
    \item отмена карточного платежа работает по-своему в~зависимости
      от~того, на~какой стадии (авторизации, клиринга или расчёта)
      находится транзакция в~этот момент.
  \end{enumerate}
}

Разделение \textit{на авторизацию и~клиринг} продиктовано скоростью
и~масштабом. Банковский перевод допускает задержку в~минуты или часы
(СБП сократила её до~секунд), и~каждый перевод обрабатывается
индивидуально. Карточный платёж должен завершиться за~секунды:
покупатель стоит у~кассы, и~продавец не может позволить себе очередь
из~десяти человек, каждый из~которых ждёт межбанковский перевод.

Разделение \textit{на авторизацию и~расчёт} даёт мгновенный ответ
без мгновенного перемещения денег: эмитент обещает заплатить, продавец
отпускает товар, фактический перевод происходит позже, пакетами,
через расчётную инфраструктуру карточной сети. Встречные потоки
взаимозачитываются, и~объём межбанковских проводок резко сокращается
по~сравнению с~индивидуальным расчётом каждой транзакции.

\enlargethispage{-5\baselineskip}
Жизненный цикл карточной транзакции, от~прикладывания карты
к~терминалу до~расчёта, подробно разбирается в~главе~\ref{ch:tx-lifecycle};
формат сообщений на~этапе авторизации рассматривается в~главе~\ref{ch:iso8583}.
Если перевод между двумя счетами в~разных банках~--- это согласованное
изменение записей в~двух местах, то карточная транзакция~--- это обещание
заплатить, которое выполняется позже. Что произойдёт, если эмитент даст
обещание, а~ночью обанкротится до~клиринга?

Расчётный риск в~этом промежутке принимает на~себя платёжная сеть: правила
Visa и~Mastercard обязывают её возместить эквайеру неурегулированную
сумму, если эмитент не~исполнит обязательство по~расчёту, а~участники
с~повышенным риском вносят обеспечение деньгами, аккредитивами или
гарантиями материнской структуры. Механизм покрытия разобран
на~примере Visa в~разделе~\ref{sec:visa-settlement-guarantee}.

Внутрироссийские карточные операции через ОПКЦ НСПК завершаются клирингом
и~расчётом в~следующий операционный день: ОПКЦ формирует клиринговую
позицию, а~расчётные услуги по~национальным платёжным инструментам
Банк России оказывает по~корреспондентским счетам участников~\cite{fz-161}.
Международные операции Visa и~Mastercard рассчитываются по~календарю
клиринговой сессии конкретной валюты~--- для долларовых платежей расчёт
обычно проходит на~T+1~\cite{visa-core-rules,mc-gcms-rm-2021}.
На~каждом календарном дне эмитент держит на~счёте уполномоченного
расчётного банка сети (см.~раздел~\ref{sec:banks-correspondent})
обеспечение под~неурегулированные авторизации~--- денежный депозит,
аккредитив или гарантию материнской структуры. Если у~эмитента
к~моменту утреннего расчётного цикла не~хватает остатка, сеть
закрывает позицию из~этого обеспечения и~параллельно инициирует
приостановку или прекращение членства эмитента в~сети.

Глава~\ref{ch:ecosystem} разбирает участников карточной экосистемы~---
эмитента, эквайера, сеть, процессор, шлюз, фасилитатор~--- и~экономику
взаиморасчётов через interchange; глава~\ref{ch:issuer-acquirer}
углубляется в~функции эмитента и~эквайера: CMS, TMS, маршрутизацию,
обработку диспутов. К~карточной транзакции как к~технической
последовательности сообщений книга возвращается в~главе~\ref{ch:iso8583}
на~уровне формата ISO~8583.

